% Tier 1 parity doc: unicode-math with an OpenType math font.
% Exercises mlist_to_hlist + OpenType math (sessions 92-108).
\documentclass[12pt]{article}
\usepackage{amsmath}
\usepackage{unicode-math}
\setmainfont{LibertinusSerif-Regular.otf}
\setmathfont{LibertinusMath-Regular.otf}
\begin{document}
\section*{Unicode math}

Inline: the value $α^2 + β^2 = γ^2$ with $∑_{i=1}^{n} i = \frac{n(n+1)}{2}$.

\begin{align}
  ∫_0^∞ e^{-x^2}\,dx &= \frac{\sqrt{π}}{2}, \\
  \nabla × \mathbf{F} &= \begin{pmatrix} ∂_y F_z - ∂_z F_y \\
                                          ∂_z F_x - ∂_x F_z \\
                                          ∂_x F_y - ∂_y F_x \end{pmatrix}.
\end{align}

A radical and a fraction over a delimiter:
\[
  \sqrt[3]{\frac{α+β}{γ}} \quad
  \left( \sum_{k=0}^{∞} \frac{x^k}{k!} \right) = e^x .
\]

\end{document}
